% Generated draft: requires mathematical and source review.
\documentclass[11pt]{article}
\usepackage[margin=25mm]{geometry}
\usepackage{fontspec}
\setmainfont{Latin Modern Roman}
\usepackage{amsmath,amssymb,mathtools,unicode-math}
\setmathfont{Latin Modern Math}
\usepackage{enumitem,needspace,xurl,hyperref}
\setlength{\parindent}{0pt}
\setlength{\parskip}{5pt}
\setlength{\emergencystretch}{3em}
\newcommand{\N}{\mathbb{N}}
\newcommand{\Z}{\mathbb{Z}}
\newcommand{\Q}{\mathbb{Q}}
\newcommand{\R}{\mathbb{R}}
\newcommand{\C}{\mathbb{C}}
\newcommand{\F}{\mathbb{F}}
\newcommand{\A}{\mathbb{A}}
\newcommand{\PP}{\mathbb{P}}
\newcommand{\E}{\mathbb{E}}
\newcommand{\eps}{\varepsilon}
\newcommand{\dd}{\,\mathrm{d}}
\begin{document}
\section*{UnsolvedMath ID: 30003024}
\textbf{Source problem code:} catalogue entry
\subsection{Definitions and mathematical statement}
Write \(\operatorname{rk}\) for Morley rank. In the standard definable interpretation, a permutation group \((X,G)\) of finite Morley rank consists of a definable group \(G\) acting definably and faithfully on a nonempty definable set \(X\), with both having finite Morley rank. The group \(G\) is connected if it has no proper definable subgroup of finite index. Its action is transitive if for every \(x,y\in X\) there exists \(g\in G\) with \(g\cdot x=y\). For a positive integer \(k\), the action is generically \(k\)-transitive if its diagonal action on \(X^k\), given by
\[
g\cdot(x_1,\ldots,x_k)=(g\cdot x_1,\ldots,g\cdot x_k),
\]
has an orbit of Morley rank \(\operatorname{rk}(X^k)=k\operatorname{rk}(X)\).

For an algebraically closed field \(K\), let \(\mathbb P^n(K)\) be the set of one-dimensional linear subspaces of \(K^{n+1}\), and let
\[
\operatorname{PGL}_{n+1}(K)=\operatorname{GL}_{n+1}(K)/\{\lambda I_{n+1}:\lambda\in K^\times\}
\]
act on \(\mathbb P^n(K)\) by its natural projective action.

\textbf{Question.} Suppose \((X,G)\) is a transitive permutation group of finite Morley rank, \(G\) is connected, \(n=\operatorname{rk}(X)\), and the action is generically \((n+2)\)-transitive. Must there exist an algebraically closed field \(K\), a group isomorphism
\[
\varphi:G\longrightarrow\operatorname{PGL}_{n+1}(K),
\]
and a bijection
\[
f:X\longrightarrow\mathbb P^n(K)
\]
such that
\[
f(g\cdot x)=\varphi(g)\cdot f(x)
\qquad(g\in G,\ x\in X)?
\]

\textbf{Unresolved qualifications.} The displayed conclusion spells out equivalence as equivalence of permutation actions. The supplied record does not specify whether the equivalence must be definable, whether \(K\) must be interpretable in the ambient structure, or which parameters are allowed. These potentially stronger requirements remain unresolved. The record supplies no explicit restriction \(n\geq1\); as the finite Morley rank of a nonempty set, \(n\) is a nonnegative integer, and no positive-rank hypothesis has been added. The ambient structure and its parameter conventions are also unspecified.
\subsection{Short English statement}
Must every transitive action of a connected group of finite Morley rank on a set of rank \(n\), with a full-rank orbit on \((n+2)\)-tuples, be equivalent to the natural projective linear action over an algebraically closed field? The required definability of this equivalence is unspecified. This is a statement-only draft requiring human mathematical review.
\subsection{Sources}
\begin{itemize}
\item Supplied local UnsolvedMath registry record, permanent ID \(30003024\), title: Generic Transitivity and Projective Linear Actions; source problem code: \texttt{catalogue entry}. Collection: Oberwolfach Reports Open Problems. The record gives proposed year 2016 and names no individual proposer.
\item Supplied report citation: \emph{Model Theory: groups, geometry, and combinatorics} (2016), Oberwolfach Reports, pp. 5--50, DOI 10.4171/owr/2016/1. Supplied source and literature link, labelled Oberwolfach Report 1/2016: \url{https://doi.org/10.4171/owr/2016/1}. The background describes these reports as records of mathematical workshops and their open-problem discussions; no precise location within the report is supplied for this question.
\item Literature assessment dated exactly \texttt{2026-08-21}: supplied dated curation, not independently verified. The record labels the problem open, but its assessment says that the exact record contains no mathematical proposition beyond a remark that a conjecture is clear, and therefore retains the open label under a triage rule. This description conflicts with the explicit mathematical question supplied here. The background requests independent verification of time-sensitive status and cited claims.
\item The supplied statement status is \texttt{corrected\_verified}; the accompanying metadata reports manual semantic, typographic, and LaTeX review for self-containment and fidelity. These are supplied review claims, not independent verification of the source or mathematical correctness.
\end{itemize}
\end{document}
